% TOP_PROBLEM: {"id":30007541,"problem_number":"LOCAL-30007541","title":"Fiscal credibility and expectations","category_id":21,"category":{"id":21,"name":"economics","display_name":"Economics","description":"Open economic theory statements, published research agendas, and proposed scoped specifications; question provenance and historical priority are qualified per record.","slug":"economics","order_index":21,"created_at":"2026-10-08T00:00:00Z"},"status":"research_question","external_url":null,"legacy_ids":[],"published":false,"statement":"Identify fiscal-announcement effects on beliefs and actual spending or financing yields independently of immediate policy changes and concurrent economic news.","economics":{"original_id":30,"field":"Fiscal policy and sovereign debt","kind":"identification","formalization_basis":"adapted_research_question","source_status":"research_agenda","priority_status":"earliest_found","priority_note":"This is the earliest explicit relevant statement verified in this audit, not a claim of original priority; older literature and earlier versions may contain antecedents.","earliest_verified_reference":{"authors":["Nicolas End","Gee Hee Hong"],"title":"Trust What You Hear: Policy Communication, Expectations, and Fiscal Credibility","year":2022,"url":"https://www.imf.org/-/media/files/publications/wp/2022/english/wpiea2022036-print-pdf.pdf","doi":"10.5089/9798400200748.001","locator":"Section 1, Introduction, paragraph contrasting monetary and fiscal expectation channels","evidence_summary":"The authors explicitly say considerably less is known about the fiscal expectations channel, then provide empirical evidence on fiscal communication and credibility. The experimental mediation target below is an adaptation."},"current_reference":{"authors":["Nicolas End","Gee Hee Hong"],"title":"Trust What You Hear: Policy Communication, Expectations, and Fiscal Credibility","year":2022,"url":"https://www.imf.org/-/media/files/publications/wp/2022/english/wpiea2022036-print-pdf.pdf","doi":"10.5089/9798400200748.001","locator":"Section 1, Introduction, paragraph contrasting monetary and fiscal expectation channels","evidence_summary":"The authors explicitly say considerably less is known about the fiscal expectations channel, then provide empirical evidence on fiscal communication and credibility. The experimental mediation target below is an adaptation."},"formalization_note":"The paper explicitly notes limited knowledge of the fiscal expectations channel and supplies partial evidence. This intervention/mediation design is adapted."},"statusQualification":"Published question or research agenda verified at the cited locator. The mathematical specification is adapted for this catalogue and includes additional assumptions; source publication does not establish first-ever priority or certify this exact model remains unsolved. The paper explicitly notes limited knowledge of the fiscal expectations channel and supplies partial evidence. This intervention/mediation design is adapted.","statusReviewedAt":"2026-10-08"}
% FIRST_POSTED: 2026-10-08
% ENTRY_KIND: statement_only
% !TeX program = lualatex
\documentclass[11pt]{article}
\usepackage{fontspec}
\usepackage[a4paper,margin=25mm]{geometry}
\usepackage{amsmath,amssymb,mathtools}
\usepackage{xurl}
\usepackage[hidelinks]{hyperref}
\newcommand{\sref}[2]{\hyperlink{source:#1:#2}{[S#2]}}
\setlength{\emergencystretch}{3em}
\begin{document}
% problemId:30007541
\section*{Fiscal credibility and expectations}\label{30007541}
\subsection{Definitions and mathematical statement}
Fix a fiscal authority’s announced future tax or spending path \(A\), holding immediate implemented policy fixed. Randomly or quasi-randomly vary credible communication \(Z_i\) among households, firms, or investors exposed to the same path. Let \(B_i(Z_i)\) record beliefs about implementation, future disposable income, and debt sustainability; let \(X_i(Z_i)\) be measured spending, investment, or required financing yield. Define the announcement intention-to-treat effect
\[
\tau_X=E[X_i(1)-X_i(0)],
\]
and, for a predetermined credibility measure \(c_i\), its conditional counterpart \(E[X_i(1)-X_i(0)\mid c_i=c]\).
Identify these effects separately from concurrent fiscal actions and news about economic fundamentals. For a belief-mediated effect, state additional exclusion and sequential-ignorability assumptions, or provide bounds when communication directly changes attention or uncertainty. Measure persistence and compare promised with implemented policy so that credibility is not retrospectively defined by favorable outcomes. To infer sovereign financing effects, specify market spillovers and whether treatment changes the information held by marginal investors. Assess whether estimates transport across institutions and fiscal regimes. The target is causal expectation and credibility effects under a declared announcement design, not the claim that fiscal communication has no known effects or that one observed spread response identifies credibility.

\par\textbf{Formulation and scope.} The paper explicitly notes limited knowledge of the fiscal expectations channel and supplies partial evidence. This intervention/mediation design is adapted.

\par\textbf{Open-question provenance.} An explicit question or research agenda was verified at the source locator. This does not by itself certify that every later version remains unresolved \sref{30007541}{1}.

\par\textbf{Historical priority.} This is the earliest explicit relevant statement verified in this audit, not a claim of original priority; older literature and earlier versions may contain antecedents.

\subsection{Short English statement}
Identify fiscal-announcement effects on beliefs and actual spending or financing yields independently of immediate policy changes and concurrent economic news.

\subsection{Sources}
\begin{itemize}\raggedright
\item[\textnormal{[S1]}]\hypertarget{source:30007541:1}{} Nicolas End, Gee Hee Hong. 2022. Trust What You Hear: Policy Communication, Expectations, and Fiscal Credibility. Section 1, Introduction, paragraph contrasting monetary and fiscal expectation channels.
\url{https://www.imf.org/-/media/files/publications/wp/2022/english/wpiea2022036-print-pdf.pdf}.
DOI: 10.5089/9798400200748.001.
{\small\textit{Earliest verified question/agenda reference; historical priority qualified below. The authors explicitly say considerably less is known about the fiscal expectations channel, then provide empirical evidence on fiscal communication and credibility. The experimental mediation target below is an adaptation.}}
\end{itemize}
\end{document}
